\section{Simulation Design}
\label{sec:simulation}

Section \ref{sec:model} defines a unit transfer benchmark and an extension stage game with occupation choice, offers and acceptances, target selection, success probabilities, and rival loot. The simulations use those same primitives in a large stochastic population where repeated matching and endogenous wealth heterogeneity make closed form characterization impractical. The section records how those theoretical objects are implemented in code and which outcomes are stored for the figures and tables.

\subsection{State variables, timing, and implementation}

The simulation state in period $t$ contains four objects: the wealth vector $w_t$, the role vector $r_t$, the patron-thief contract network, and the guard network. Role values are
\[
\{\text{honest}, \text{stay}, \text{freelance thief}, \text{hired thief}, \text{patron}, \text{guard}, \text{elite}\},
\]
with the understanding that an elite can also be a patron or a protected household once the relevant threshold is crossed. After wealth is observed, agents are classified into rich and elite sets by threshold rules based on current mean wealth, and the code then reproduces the extension stage game in the same order as Section \ref{sec:model}. Rich households choose among staying, self raiding, and making theft offers. Elite households additionally compare continued patronage with guard hiring through \eqref{eq:guard-threshold}. Poor households accept at most one offer or choose freelance predation or withdrawal. Active thieves then select targets, success is governed by \eqref{eq:success-probability}, contractual payments and guard wages are applied, punishment is realized where relevant, and the next state is recorded.

Deterministic code handles the exact event accounting required by the first honest night, while stochastic code governs the richer transition once matching, targeting, and enforcement interact. The baseline seed is fixed at 613, every sweep uses reproducible seed offsets tied to the parameter grid, and every displayed result is generated from the same pipeline without hand-edited intermediates.

Event time is aligned on the first honest night, so the initial perturbation remains distinct from later adaptation. The baseline routine isolates the direct consequence of one honest household staying home. The later stochastic routines then measure what happens once wealth heterogeneity changes occupational choice, target selection, and protection decisions. Parameter changes can therefore be read through the channel they affect: amplification of the local asymmetry, the profitability of delegation, or the private value of guarding.

\subsection{Mapping the theory into simulation rules}

The benchmark simulation stays as close as possible to Theorem \ref{thm:no-stationary-one-honest}. Agents begin on a balanced cycle in the unit transfer environment, one agent is permanently honest, and on the first night the honest agent stays home while the old cycle is otherwise followed exactly. This reproduces the accounting in Lemma \ref{lem:first-night}. The same collision rule used in the benchmark game is preserved, although no congestion occurs along the initial cycle because each target is chosen once. In later benchmark periods, agents compare their wealth to the cutoff in \eqref{eq:threshold}. Households below the threshold continue personal raiding along the cycle, while households above it stay home. Contracting, guarding, and stochastic targeting are absent at this stage, so any departure from the balanced cycle is driven directly by the honest perturbation and the wealth-dependent exposure term.

Once proportional loot is introduced, role thresholds are defined relative to current mean wealth. An agent becomes rich when his wealth exceeds $\theta_R \bar{w}_t$, where $\bar{w}_t$ is the cross-sectional mean, and becomes elite when wealth exceeds the higher threshold $\theta_E \bar{w}_t$. These thresholds are reduced-form counterparts of the occupational comparisons in Section \ref{sec:model}. Rich households are those for whom exposure makes personal raiding relatively unattractive, while elites are those for whom the stock of wealth at risk makes guarding salient. In the baseline calibration, the thresholds preserve event-time separation between the first honest shock and later institution building while still allowing the patronage regime to emerge within the simulation horizon.

Rich households are sorted by wealth, and candidate thieves are drawn from the poorest available agents who are neither permanently honest nor already employed. Sorting employers and matching them to the poorest available workers is the computational analogue of the offer stage in the extension game. A contract is formed only if the feasibility condition in Lemma \ref{lem:contract-feasibility} holds with current state variables. When a contract is feasible, the simulation applies the bargaining rule in \eqref{eq:bargaining-rule}: the thief receives his outside option plus the share of residual surplus not captured by the patron, and the patron receives the rest. The hired thief then targets households in the high value poor set $H_t$, while freelancers target vulnerable households by expected loot. Rival loot remains one object here as well. When several freelancers or hired thieves converge on the same house, the code applies the same congestion logic as \eqref{eq:success-probability} and \eqref{eq:contract-update}: attackers draw from one vulnerable stock, so each additional attacker lowers the expected return of the others.

The target rule keeps the direction of predation explicit. In the patronage regime, poor thieves usually rob poor households for rich patrons. Physical theft remains largely horizontal among the nonelite, while the net transfer is vertical because part of the loot moves upward through the contract. The candidate target set excludes the patron and currently protected elite households. The patron's house remains occupied in the patronage phase, and betrayal risk enters only through the residual elite vulnerability summarized by $\chi_t$ once patronage begins to unravel.

Elites then face a second occupational comparison, namely whether the expected surplus from another patron-thief contract is smaller than the expected value of avoided elite loss from a guard. If inequality has risen and the prey base among the poor has thinned, the guarding condition in \eqref{eq:guard-threshold} is satisfied for some elite households. Guards are again recruited from the poorest available agents, but only if the guard wage exceeds their best alternative. Once hired, a guard reduces the success probability of raids on the protected house and introduces a punishment probability if an attempted elite raid fails. This gives the simulation an explicit mechanism by which formal protection suppresses open predation without removing the underlying population willing to steal.

Across the simulation, chronology runs through three stages that match the theory. The first stage is the deterministic first-night perturbation in the unit transfer benchmark. The second is the patronage stage in the proportional environment, where class differentiation and delegated predation emerge. The third is the enforcement stage, where some elites replace delegated predation with guards and punishment. Endogenous threshold crossings generate that chronology, so calendar dates do not drive the switch and transition timing remains an outcome of wealth dispersion, target depletion, outside options, and enforcement effectiveness.

\subsection{Baseline calibration}

Baseline parameter values are chosen in the internal units of the model and are not estimated from a specific dataset. The number of agents is set to $N=60$, which is large enough to generate meaningful cross-sectional dynamics and small enough to keep network visualizations interpretable. Initial wealth is normalized to one for every agent. The proportional theft intensity is $\gamma=0.35$, so a successful raid extracts a substantial but not total fraction of victim wealth. The rich threshold is $\theta_R = 1.15$ times current mean wealth, while the elite threshold is $\theta_E = 2.10$ times current mean wealth. These values allow early windfalls to matter without producing an elite stratum immediately.

The remaining parameters pin down the costs and comparative advantages of the occupational choices. The cost of theft effort is $c_T = 0.02$. The contracting cost is $k = 0.03$. The baseline guard wage is $g = 0.08$. The target thief share under a contract is $\alpha = 0.45$, while the fixed wage component is calibrated at $s = 0.03$. The success probability of raiding an occupied but unguarded household is low but positive, and the success probability against a guarded elite household is lower still. Punishment after a failed guarded raid is likely but not certain. In words, the calibration assumes a society in which unprotected wealth is vulnerable, protection is costly, and the poor can be induced to work as thieves or guards because their outside options are weak. Table \ref{tab:parameters} reports the core baseline values used in the main figures.

\begin{table}[tbp]
    \centering
    \caption{Baseline simulation parameters. The values define the benchmark environment used in the main figures, while the parameter sweeps vary the key margins systematically around this baseline.}
    \label{tab:parameters}
    \begin{tabular}{lll}
\toprule
Symbol & Value & Description \\
\midrule
$N$ & 60 & Number of agents \\
$\gamma$ & 0.350 & Proportional theft intensity \\
$\bar{w}$ & 1 & Initial wealth per agent \\
$\theta_R$ & 1.150 & Rich threshold multiplier \\
$\theta_E$ & 2.100 & Elite threshold multiplier \\
$s$ & 0.030 & Base contract wage \\
$\alpha$ & 0.450 & Thief loot share under contract \\
$\kappa$ & 0.030 & Contracting cost \\
$g$ & 0.080 & Guard wage \\
$q$ & 0.980 & Guard effectiveness \\
\bottomrule
\end{tabular}

\end{table}

Parameters that govern the first night logic mirror the benchmark theory. Parameters that govern patronage and enforcement are chosen so that each regime can emerge for reasons consistent with the propositions in Section \ref{sec:model}. The baseline outside option of poor thieves is low enough that patronage can be strictly profitable, but not so low that every rich agent hires immediately. Guard effectiveness is high enough that protection can dominate once wealth is sufficiently concentrated, but not so high that guards appear before a meaningful patronage phase.

The baseline path is deliberately moderate: inequality rises without exploding instantly, guards appear only after a visible patronage phase, and collapse remains possible without becoming inevitable. A more aggressive calibration would produce sharper pictures, but it would also obscure whether the mechanisms emerge from the stated thresholds or from extreme parameter values.

\subsection{Parameter sweeps and regime classification}

Systematic parameter sweeps then vary the margins highlighted in the theory. Theft intensity $\gamma$ is varied to measure how sensitive class formation is to the fraction of wealth that can be appropriated in a single raid. Exposure risk is varied through the success probability of attacks on occupied but unguarded households, which determines how costly it is for a rich household to rely on occupancy alone. Contracting cost $k$ and patron bargaining power $\beta$ are varied to map the feasibility and profitability of delegation. Guard effectiveness and punishment severity are varied to locate the transition from patronage to protection. The outside option of poor thieves is varied through the congestion parameter $\lambda_t$ and the availability of guard employment. Finally, the size of the initial honest set is varied so that the simulations can distinguish a single local perturbation from a broader coordinated honesty shock.

Each period in each simulation path is classified into one of six regimes: balanced cycle, transition, patronage, guard regime, coexistence, or collapse. The balanced-cycle label is reserved for periods in which almost all agents continue the reciprocal pattern and inequality remains negligible. Transition denotes periods after the first honest disturbance but before patronage or guarding has become dominant. Patronage denotes periods with active patron-thief contracts and no meaningful guard presence. Guard regime denotes periods in which protected elites and guards are established at scale. Coexistence denotes periods in which patronage and guarding coexist in nontrivial numbers. Collapse denotes periods in which the poor prey base is so depleted that predatory and protective roles no longer reproduce the earlier structure. These labels are not rhetorical. They are deterministic functions of role counts, inequality thresholds, and guard presence stored in the simulation output. The phase diagram in Section \ref{sec:results} is built directly from these classifications.

Fixed coding rules determine regime classification. Regimes are not inferred from unsupervised clustering, and paths are not relabeled ex post to fit the figures. The same classification rule is imposed across the entire parameter grid, which is necessary if a phase diagram is to carry comparative meaning.

\subsection{Recorded outcomes}

For each path the code stores agent level wealth, role histories, employer-worker links, target links, guard assignments, realized and expected welfare components, inequality measures, and first entry times into patronage and guard states. These objects support the figures that the theory alone cannot provide: event time wealth trajectories by eventual role, distributional heatmaps and quantile ribbons, Lorenz curves at key dates, phase diagrams across parameter pairs, and network snapshots of theft, patronage, and guarding. The same stored objects support the tabulated summaries of first patron entry, first guard entry, final inequality, and dominant long run regime.

Recorded outcomes answer a limited set of analytical questions tied to the theory: how the first honest perturbation propagates through the wealth distribution, when delegated predation becomes organizationally stable, and when organized protection replaces or disciplines predation. Event time plots establish the timeline. Distributional plots establish stratification. Comparative statics show that patronage and guarding do not depend on one baseline path. Welfare decompositions track what is lost and what, at times, is partially recovered as regimes change.
